% Auto-generated by scripts/render_lanl_replication_results.py from frozen public LANL evidence.
% Paired CSV SHA-256: 6dfe2e18ed65d2f54d6a73ce5cb8b9848b0ecdf78e3fd36bbf469a52396f14db
% Control manifest SHA-256: 190190ac1b75495fdfb18422da8b727bdac88553852edba8d236fb3d544957ff
\subsection{Frozen Second-Cohort Observations and Missing-Control Sensitivity}
The preregistered LANL second cohort (original red-team rows 26--50)
contains 25 targets, of which 18 have eligible temporal controls and
7 remain unmatched. Controls are \emph{unlabeled}, not verified benign.
Table~\ref{tab:lanl-second-frozen} summarizes three prespecified process
features. The median of within-pair differences need not equal the
difference between target and control marginal medians.
\begin{table}[ht]
\centering
\caption{Frozen LANL second-cohort descriptive comparisons.}
\label{tab:lanl-second-frozen}
\begin{tabular}{lrrrr}
\toprule
Feature & T/C medians & Paired $\Delta$ & T$>$C & Bounds \\
\midrule
Process events & 64/9 & 44.5 & 16/18 & [0.64,0.92] \\
Unique processes & 27/3.5 & 16.5 & 17/18 & [0.68,0.96] \\
Process starts & 49/5.5 & 29 & 16/18 & [0.64,0.92] \\
\bottomrule
\end{tabular}
\end{table}
The bounds apply only to the 25 selected targets: unmatched controls
have unknown nonnegative process counts. The lower bound treats all
unmatched comparisons as not target-higher; the upper bound counts
an unmatched comparison as potentially target-higher only when its
target count is positive. These are logical finite-cohort bounds,
\emph{not} confidence intervals or population estimates.
A post-hoc audit found 4 distinct users and 2 source computers among
matched targets and 13 overlapping-window episodes in the frozen cohort.
Repeated users, source hosts and overlapping windows preclude treating
pair-level sign-test $p$ values or Wilson intervals as calibrated
population-level inference \cite{cameron2015cluster}.
The three process features are correlated and do not represent
independent discoveries. The data do not establish attack-detection
accuracy, verified benign controls, causal effects, or the effectiveness
of the proposed eight-factor device-trust access model. Independent
endpoint-posture measurements and adjudicated access outcomes remain
necessary before external effectiveness claims.
